

\documentclass[12pt,letterpaper]{article}
\usepackage[utf8]{inputenc}
\usepackage[T1]{fontenc}
\usepackage[english]{babel}
\usepackage{amsmath}
\usepackage{amsfonts}
\usepackage{amssymb}
\usepackage{graphicx}
\usepackage{booktabs}
\usepackage{geometry}
\usepackage{hyperref}
\usepackage{float}
\usepackage{placeins}
\usepackage{microtype}
\usepackage{caption}
\captionsetup{font=small,labelfont=bf}

% ---- add to the preamble of AMLmain (after \usepackage{hyperref} is fine) ----
\usepackage{xcolor}
\usepackage{enumitem}
% placeholder for text to be written after the experiments; renders in red so nothing is forgotten
\newcommand{\fillin}[1]{\textcolor{red}{[#1]}}
% numbers produced by summarize.py (FinalTestAUC, GeomTestAUC, RefPatValAUC, ...)
\input{results/tables/macros.tex}



\geometry{top=2.5cm, bottom=2.5cm, left=2.5cm, right=2.5cm}

\title{\vspace{5 cm}Deep Learning for Medical Image Classification\vspace{1cm}}
\author{Edoardo Buttazzi}

\date{July 2026\\\vspace{1.5cm} \textbf{Applied Machine Learning-Advanced }}
\begin{document}
\maketitle
\clearpage
\tableofcontents
\newpage
\section{Introduction and Project Objective}
This project investigates binary classification of pediatric chest radiographs into \textbf{NORMAL} and \textbf{PNEUMONIA} using transfer learning implemented in PyTorch. Beyond reporting discrimination on the source test set, it asks which design choices actually matter and whether the network does better than technical shortcuts present in the data.

The analysis proceeds in four steps: (i) a quantitative audit of the dataset, which reveals that native image geometry alone separates the classes; (ii) a patient-grouped validation protocol in which model, checkpoint and decision threshold are selected on validation only; (iii) a comparison against classical and deep baselines and one-factor-at-a-time ablations of architecture, imbalance strategy, padding and learning rates; (iv) error analysis and Grad-CAM with a padding ablation measured on the whole test set. All code, configurations and generated tables are part of the accompanying repository.
\newpage

\section{Architectural Selection}
\textbf{DenseNet121} was selected as the backbone for distinguishing NORMAL and PNEUMONIA chest radiographs because of its established role in advanced medical-image classification. Its contribution to the state of the art in chest-radiograph analysis is illustrated by \textbf{CheXNet}, a system based on a 121-layer densely connected network. In 2017, the CheXNet study reported state-of-the-art results across the 14 thoracic disease categories of ChestX-ray14, including pneumonia~\cite{rajpurkar2017chexnet}. This evidence provides a relevant literature-based justification for choosing DenseNet121 for the present task.

The architecture's main strength is its \textbf{dense connectivity}. Within each dense block, every layer receives the feature maps produced by all preceding layers. This design encourages feature reuse, improves the propagation of information and gradients, and reduces the number of parameters needed to learn useful representations~\cite{huang2017densenet}. For chest radiographs, access to features extracted at different depths is valuable because classification can require combining fine texture and contrast patterns with broader anatomical structures.

DenseNet121 also supports \textbf{transfer learning through ImageNet-pretrained weights}. This provides an initialization for learning visual representations from a development pool of limited size, while its parameter efficiency makes fine-tuning practical within the available computational resources. The pretrained backbone can then be adapted to the binary task through a dedicated classification head and progressive unfreezing, as described in the following sections.

The choice therefore rests on demonstrated relevance to chest-radiograph research, efficient feature reuse, and compatibility with the project's transfer-learning strategy. The state-of-the-art reference concerns the published CheXNet benchmark results at the time of that study; performance on the present pediatric dataset is assessed through the project's own evaluation.

\section{Dataset and Exploratory Data Analysis}
\label{sec:dataset}

\subsection{Source, population, and prediction target}
The project uses the chest-radiograph component of the public dataset released by Kermany et al.\ \cite{kermany2018,kermanydata}. The source study describes anteroposterior chest X-rays collected during routine care at Guangzhou Women and Children's Medical Center, China, from children aged one to five years. Image quality was screened before inclusion; two physicians assessed the diagnostic labels, and a third reviewed the evaluation set. These statements describe the source collection rather than an independent clinical re-annotation in this project.

Each observation is an image, with \textbf{NORMAL} encoded as $y=0$ and \textbf{PNEUMONIA} as $y=1$. Bacterial and viral source labels are merged for the binary task. They remain available in filenames for descriptive analysis, but filenames, directory names, and subtype strings are never valid predictive inputs. The dataset supplies image-level diagnostic categories, without lesion masks, bounding boxes, or complete patient-level demographic and acquisition metadata. The task therefore concerns classification against these labels; it does not establish the cause of infection or clinical suitability for an individual patient.

\subsection{Archive inventory and split composition}
The inventory is computed directly from \texttt{ChestXRay2017.zip}. Only image members under the source class directories are included. Hidden files and macOS resource-fork entries are excluded, so archive-entry counts are not confused with sample counts. All \textbf{5,856 images} decode successfully. The archive contains a development directory named \texttt{train} and a separate \texttt{test} directory; it does not contain a validation directory. Table~\ref{tab:source-counts} reports the verified counts before deduplication or augmentation.

\begin{table}[htbp]
\centering
\caption{Verified composition of the local source archive. The bacterial and viral columns partition the PNEUMONIA class. Counts refer to images, not verified unique patients.}
\label{tab:source-counts}
\small
\begin{tabular}{lrrrrr}
\toprule
Source subset & NORMAL & Bacterial & Viral & Total & Pneumonia \\
\midrule
Development pool & 1,349 & 2,538 & 1,345 & 5,232 & 74.22\% \\
Source test & 234 & 242 & 148 & 624 & 62.50\% \\
\midrule
Entire archive & 1,583 & 2,780 & 1,493 & 5,856 & 72.97\% \\
\bottomrule
\end{tabular}
\end{table}

The original experimental notebook constructs the fitting and validation subsets with \texttt{train\_test\_split}, using \texttt{test\_size=0.2}, \texttt{random\_state=42}, and stratification by the binary target. The fitting subset contains \textbf{4,185 images}: 1,079 NORMAL and 3,106 PNEUMONIA. Validation contains \textbf{1,047 images}: 270 NORMAL and 777 PNEUMONIA. The source test subset retains its 624 images. These are historical image-level counts, retained to make the saved training results traceable. The split seed applies to the original partitioning operation. Augmentation is applied after partitioning and exclusively during fitting.

Figure~\ref{fig:class-balance} shows a \textbf{11.72 percentage-point difference} in pneumonia prevalence between the development pool and source test set. The development imbalance is approximately $2.88{:}1$, whereas the test imbalance is $1.67{:}1$. A classifier predicting PNEUMONIA for every image would obtain 74.22\% development accuracy and 62.50\% test accuracy, with zero specificity and 50\% balanced accuracy. These are arithmetic reference values, not fitted-model results. The positive prevalence also sets the reference precision for random ranking: approximately 0.742 in development and 0.625 in the source test set. Accuracy, precision, and precision--recall summaries should therefore be interpreted alongside class counts and sensitivity/specificity.

\begin{figure}[htbp]
\centering
\includegraphics[width=\linewidth]{images/class_balance.png}
\caption{Observed source class balance, before any resampling. Bar lengths show proportions; labels report absolute counts.}
\label{fig:class-balance}
\end{figure}

\subsection{Image dimensions, intensity, and technical variability}
Exploratory measurements are computed on the unaugmented images. Native widths range from 384 to 2,916 pixels, heights from 127 to 2,713 pixels, and width-to-height ratios from 0.835 to 3.379. The archive contains 5,573 single-channel images and 283 RGB-encoded images. All RGB-encoded records occur in the development PNEUMONIA class. This encoding difference is a potential acquisition or export confounder; it does not by itself prove that color information differs or that the model uses it. A common grayscale conversion followed, when required, by replication into three channels provides a consistent representation for inspection and model input.

Table~\ref{tab:image-statistics} quantifies the variation across source subset and class. Mean intensity $\mu$ and contrast $\sigma$ are measured after grayscale conversion and Lanczos resizing to $256\times256$, exclusively for computing comparable whole-frame descriptors. Intensities are divided by 255. Native dimensions and aspect ratios are measured before this resizing. For $M$ grayscale pixels $g_j\in[0,1]$,
\[
\mu=\frac{1}{M}\sum_{j=1}^{M}g_j,
\qquad
\sigma=\sqrt{\frac{1}{M}\sum_{j=1}^{M}(g_j-\mu)^2}.
\]
These are image statistics, not calibrated measures of tissue density. They include borders, background, and anatomy outside the lungs, and depend on the stated resampling procedure.

\begin{table}[htbp]
\centering
\caption{Observed image properties by source subset and class. Width and height are medians with interquartile ranges; the remaining columns report medians. N = NORMAL; P = PNEUMONIA.}
\label{tab:image-statistics}
\small
\setlength{\tabcolsep}{4pt}
\begin{tabular}{lrrrrr}
\toprule
Subset/class & Width, px [IQR] & Height, px [IQR] & $W/H$ & $\mu$ & $\sigma$ \\
\midrule
Development N & 1,638 [1,466--1,820] & 1,325 [1,152--1,541] & 1.223 & 0.479 & 0.241 \\
Development P & 1,168 [1,000--1,368] & 776 [640--968] & 1.488 & 0.483 & 0.214 \\
Test N & 1,762 [1,533--2,030] & 1,318 [1,040--1,639] & 1.353 & 0.487 & 0.238 \\
Test P & 1,111 [992--1,256] & 736 [632--872] & 1.506 & 0.469 & 0.215 \\
\bottomrule
\end{tabular}
\end{table}

The clearest difference is geometric: NORMAL images are typically larger and less horizontally elongated than PNEUMONIA images. In the development pool, median intensity is similar across classes, while median contrast differs. Neither observation is sufficient to establish a disease mechanism. Differences may reflect acquisition, cropping, export conventions, or patient positioning as well as image content. Resizing every input to the same tensor dimensions does not necessarily remove these cues, because field of view, interpolation texture, and padding can preserve information about the original geometry.

As a descriptive check of this confounding risk, the score $s=-W$, using native image width alone, yields an empirical rank AUC of \textbf{0.8797} in the development pool and \textbf{0.9500} in the source test set. This statistic is calculated from all positive--negative image pairs, giving half credit to tied widths. No neural network, threshold optimization, or image content is involved. It demonstrates substantial class separation in a technical property and provides a concrete reason to investigate very high early-epoch metrics. It does not identify which cues DenseNet uses, and the retrospective test inspection is not an independent evaluation of a new classifier.


\subsection{Representative images and qualitative inspection}
Figure~\ref{fig:dataset-samples} shows two actual development images for each source category: NORMAL, bacterial PNEUMONIA, and viral PNEUMONIA. Selection is deterministic: within each category, images nearest the 25th and 75th percentiles of mean intensity are selected, with source-path order breaking ties. This rule illustrates technical variability without choosing examples according to model confidence or apparent ease of diagnosis. Original aspect ratios are preserved for display, and no contrast enhancement is applied.

The examples support inspection of exposure, coverage of the thorax, orientation, visible annotations, and non-lung structures. Source pneumonia labels describe the supplied category; the figure does not provide expert lesion localization. Opacification can be a relevant radiographic finding, but whole-image brightness or low contrast alone should not be interpreted as pneumonia. Likewise, a visible line or bright object should not be assigned a device type or clinical meaning without expert review. These distinctions prevent qualitative illustrations from being mistaken for verified diagnostic explanations.

\begin{figure}[p]
\centering
\includegraphics[width=0.91\linewidth,height=0.77\textheight,keepaspectratio]{images/representative_images.png}
\caption{Six source images selected by a fixed intensity-quantile rule. Rows show NORMAL, bacterial PNEUMONIA, and viral PNEUMONIA; columns show examples near the lower and upper intensity quartiles. Filenames provide traceability. Source: Kermany et al.\ \cite{kermanydata}, CC BY 4.0. Grayscale conversion and display resizing only; no diagnostic annotation or model explanation is implied.}
\label{fig:dataset-samples}
\end{figure}

\subsection{Duplicate audit and limits of patient separation}
The integrity audit compares SHA-256 hashes of compressed image bytes and of decoded grayscale pixels with their native dimensions. Both checks identify the same \textbf{30 duplicate groups}, containing 62 records and therefore 32 redundant copies. There are 24 groups within development and six within the source test set. No exact duplicate group spans the two source directories, and no group has conflicting binary labels.

These results do not certify independence of the training and validation subsets. A random image-level split of the development pool could place identical copies in different subsets. Correlated images or repeated examinations can also cause leakage without being byte-identical. This is why the revised protocol removes exact duplicates and groups the validation split by patient key (Section~\ref{sec:protocol}); experiment Q0 measures how much the image-level split of the first submission inflated validation scores.

Filename inspection provides an additional, limited check. A coarse token such as \texttt{person123} is reused across bacterial and viral namespaces: 170 such tokens appear in both source subsets. Retaining the subtype in the grouping key produces 4,097 filename groups across the archive, including 726 groups with multiple images, with no group spanning development and test. The normal-image key retains the \texttt{IM} or \texttt{NORMAL2-IM} prefix and its numeric identifier. These are naming-pattern findings, not verified patient counts. The absence of shared keys cannot replace a documented patient identifier, and reused coarse tokens cannot be taken as proof of patient overlap.

An exploratory 64-bit difference-hash screen at Hamming distance at most five returns 461,633 candidate image pairs, including 103,449 pairs across the source directories. These are candidate counts only: similar chest geometry makes this simple screen highly non-specific. They are not counted as duplicates and are not used for automatic exclusion. The exact-hash findings above are the confirmed duplication result.

\subsection{Bias, confounders, and reproducibility}
The audit supports several limits on interpretation. The single-center pediatric source restricts the represented population; performance on adults, other hospitals, or different imaging systems requires separate evidence. Source quality screening can reduce the representation of difficult or unreadable examinations. The binary label merges bacterial and viral cases and does not describe the full range of alternative chest abnormalities. The quantitative analysis concerns image-level labels, geometry, encoding, and intensity. These measurements define the scope of the quantitative dataset conclusions.

Class-associated resolution, aspect ratio, and encoding are measurable potential shortcuts. Representative images and descriptor plots characterize acquisition and appearance. Their scope is distinct from neural attention or expert-validated lesion localization.

The accompanying audit files record relative archive paths, source labels and subtypes, image dimensions and mode, whole-frame descriptors, hashes, and filename-group keys. The analysis is deterministic and applies no augmentation or label-dependent image processing. The source archive is read without alteration. Its distribution is attributed to Kermany et al.\ \cite{kermanydata} under CC BY 4.0. The inventory and the configuration files in the repository document the dataset analysis and every training run. Source-directory separation, exact-duplicate checks, and verified patient-level independence are distinct properties.

\FloatBarrier

\section{Methods}
\label{sec:methods}

\subsection{Pipeline}
Images are decoded in grayscale, replicated to three channels, resized so that the
longest side equals the input size and padded to a square. Training applies random
affine transformations, brightness/contrast perturbation and CLAHE; validation and
test share one deterministic transform. DenseNet121 is initialised with ImageNet
weights and its classifier replaced by dropout (0.4) and a linear layer with one
logit. The backbone is frozen for the first three epochs and then fine-tuned with
three AdamW parameter groups (early backbone, late backbone, head). All settings are
stored in \texttt{configs/default.yaml}; Table~\ref{tab:config} lists the reference
configuration and the alternatives explored in Section~\ref{sec:comparison}.

\begin{table}[htbp]
\centering
\caption{Reference configuration and alternatives evaluated in the screening experiments.}
\label{tab:config}
\small
\begin{tabular}{lll}
\toprule
Component & Reference & Alternatives (one at a time) \\
\midrule
Backbone & DenseNet121, ImageNet & ResNet50; small CNN from scratch \\
Fine-tuning & 3 frozen epochs, then all layers & linear probe; no frozen phase \\
Learning rates & $3{\times}10^{-6}$ / $7{\times}10^{-6}$ / $4{\times}10^{-4}$ & backbone $\times5$; head $/4$ \\
Loss & focal, $\gamma=2$, no class weight & BCE; BCE + \texttt{pos\_weight}; focal + $\alpha_t$ \\
Sampling & $\propto 1/\sqrt{n_c}$ & uniform; $\propto 1/n_c$ \\
Padding & reflection & constant (mean); anisotropic resize \\
Validation split & patient-grouped & image-level (previous submission) \\
\bottomrule
\end{tabular}
\end{table}

\subsection{Class imbalance: an explicit choice}
\label{sec:imbalance}
The development pool contains about 2.9 PNEUMONIA images per NORMAL image. The previous
submission combined three mechanisms whose roles were not separated: a $1/\sqrt{n_c}$
sampler, a focal loss whose coefficient $\alpha=0.85$ multiplied \emph{both} classes
(and therefore did not re-weight them), and a \texttt{pos\_weight} that was computed but
never used. In the revised code each mechanism is a configuration option and is tested in
isolation (group Q2):
\begin{align*}
\mathcal{L}_{\mathrm{BCE}}^{w} &= -\tfrac{1}{B}\textstyle\sum_i \big[w_{+}\,y_i\log p_i + (1-y_i)\log(1-p_i)\big],
  \qquad w_{+}=n_0/n_1,\\
\mathcal{L}_{\mathrm{focal}} &= \tfrac{1}{B}\textstyle\sum_i \alpha_{t,i}\,(1-p_{t,i})^{\gamma}\,\ell_i,
  \qquad \alpha_{t,i}=\alpha y_i+(1-\alpha)(1-y_i),
\end{align*}
with $\ell_i$ the per-example cross-entropy and $p_{t,i}$ the probability of the true class.
Setting $\alpha=n_0/(n_0+n_1)\approx0.26$ gives the minority class NORMAL about three times
the weight of PNEUMONIA. Re-weighting is never combined with re-sampling, so that no
correction is applied twice. \fillin{State which strategy was retained and why, citing
Table~\ref{tab:screening}.}

\section{Experimental Protocol}
\label{sec:protocol}
\textbf{Patient-grouped validation.} The archive provides no patient identifier. A
conservative key is derived from filenames (\texttt{personN} for PNEUMONIA, merging its
bacterial and viral files; \texttt{IM-NNNN} and \texttt{NORMAL2-IM-NNNN} for NORMAL).
Merging can only enlarge groups, so the resulting split is stricter, not leakier. Exact
duplicates (identical SHA-256) are removed, and validation is one fold of a
\texttt{StratifiedGroupKFold} with five folds (Table~\ref{tab:splits}). The image-level split
of the previous submission is kept as experiment \texttt{ref\_image\_split} to measure how
much leakage inflated validation.

\begin{table}[htbp]
\centering
\caption{Split composition after duplicate removal. No patient key is shared between training and validation.}
\label{tab:splits}
\small
\input{results/tables/tab_splits.tex}
\end{table}

\textbf{Selection rules fixed in advance.} (i) The checkpoint of each run is the epoch with
the highest validation ROC-AUC. (ii) The final configuration is the one with the highest mean
validation ROC-AUC in the screening. (iii) The decision threshold is the highest value that
reaches a validation sensitivity of at least 0.95, a choice that prioritises missed pneumonia
over false alarms; it is frozen before the test set is used. The test set enters only the
final evaluation.

\textbf{Budget and variability.} Screening runs use $384\times384$ inputs, eight epochs and one
seed; the selected configuration and the reference are then re-trained at $512\times512$ with
three seeds, and results are reported as mean $\pm$ standard deviation. Test ROC-AUC is reported
with a 95\% stratified bootstrap interval; the difference from the geometry-only baseline uses a
paired bootstrap on the same test images.

\section{Experimental Comparison}
\label{sec:comparison}

\subsection{Baselines}
Three non-deep baselines share the split, threshold rule and metrics of the neural models:
a constant predictor (AUC 0.5 by construction), a logistic regression on native
width, height and aspect ratio only, and a logistic regression on $64\times64$ thumbnails.
The geometry-only model reaches validation ROC-AUC \GeomValAUC{} and test ROC-AUC
\GeomTestAUC{}. A neural network must therefore be judged against this value rather than
against 0.5. Deep baselines (small CNN from scratch, DenseNet121 linear probe, ResNet50) are
part of the screening.

\subsection{Screening results}
Table~\ref{tab:screening} and Figure~\ref{fig:screening} summarise all configurations on the
validation set. Each row changes one factor relative to the reference.

\begin{table}[htbp]
\centering
\caption{Screening on the patient-grouped validation set ($384\times384$, 8 epochs). AUC: ROC-AUC;
AP: average precision; sensitivity, specificity and balanced accuracy at the validation-selected
threshold. Bold: configuration selected for the final stage. The test set is not used here.}
\label{tab:screening}
\resizebox{\linewidth}{!}{\input{results/tables/tab_screening.tex}}
\end{table}

\begin{figure}[htbp]
\centering
\includegraphics[width=0.85\linewidth]{results/figures/screening_val_auc.pdf}
\caption{Validation ROC-AUC per configuration, coloured by research question. Error bars,
where present, show the standard deviation over seeds.}
\label{fig:screening}
\end{figure}

\paragraph{Q0 -- leakage.} With the image-level split the reference configuration obtains
validation ROC-AUC \RefImgValAUC{}; with the patient-grouped split it obtains \RefPatValAUC{}.
\fillin{Interpret the difference: how much of the near-perfect early validation scores of the
first submission is explained by correlated images across the split.}

\paragraph{Q1 -- architecture and transfer learning.} \fillin{Compare scratch CNN, linear
probe, full fine-tuning and ResNet50. Does fine-tuning add over frozen ImageNet features? Does
DenseNet121 justify its choice over ResNet50 on this dataset?}

\paragraph{Q2 -- imbalance.} \fillin{Compare plain BCE, \texttt{pos\_weight}, re-sampling and
class-weighted focal loss. Expected pattern: similar ROC-AUC (ranking is largely unaffected by
re-weighting) but different operating points; the frozen-threshold rule then removes most of
the difference.}

\paragraph{Q3 -- padding.} Grad-CAM in Section~\ref{sec:xai} shows positive attribution
in reflected borders. \fillin{Report whether constant padding or anisotropic resizing changes
validation performance, and whether it reduces the attribution in padding.}

\paragraph{Q4 -- learning rates.} \fillin{Sensitivity of the result to the learning rates.}

\section{Final Model: Test Evaluation}
\label{sec:final}
The selected configuration is \textit{\FinalName}. Over \FinalSeeds{} seeds it reaches test
ROC-AUC \FinalTestAUC{}. At the validation-selected threshold ($t=\FinalThreshold$) the test
sensitivity is \FinalTestSens{} and specificity \FinalTestSpec{} (Table~\ref{tab:final},
Figures~\ref{fig:learning}--\ref{fig:confusion}).

\begin{table}[htbp]
\centering
\caption{Test results at the threshold selected on validation. Neural models: mean $\pm$ sd
over seeds; the CI is the mean of per-seed bootstrap intervals. $\Delta$AUC: paired bootstrap
difference from the geometry-only baseline, 95\% interval.}
\label{tab:final}
\resizebox{\linewidth}{!}{\input{results/tables/tab_final_test.tex}}
\end{table}

\begin{figure}[htbp]
\centering
\includegraphics[width=\linewidth]{results/figures/learning_curves_final.pdf}
\caption{Learning curves of the final runs (mean over seeds; band: min--max). Dashed:
training loss on augmented, re-sampled data; solid: validation. The dotted vertical line
marks the start of backbone fine-tuning.}
\label{fig:learning}
\end{figure}

\begin{figure}[htbp]
\centering
\includegraphics[width=\linewidth]{results/figures/roc_pr_test.pdf}
\caption{Test ROC and precision--recall curves of the final model (first seed) and the
geometry and pixel baselines. Dotted lines: chance (ROC) and prevalence (PR).}
\label{fig:rocpr}
\end{figure}

\begin{figure}[htbp]
\centering
\includegraphics[width=0.48\linewidth]{results/figures/confusion_test.pdf}
\caption{Test confusion matrix of the final model at the frozen threshold; percentages are
row-normalised (sensitivity and specificity on the diagonal).}
\label{fig:confusion}
\end{figure}

\fillin{Two or three sentences: does the network significantly exceed the geometry baseline
($\Delta$AUC interval excluding zero)? Which error type dominates and what would it mean in a
screening setting?}

\section{Error Analysis and Explainability}
\label{sec:xai}
\subsection{Purpose and class-specific attribution}
Native image dimensions can separate the source labels, but this observation does not show where the classifier obtains evidence. Grad-CAM~\cite{selvaraju2017gradcam} adds a spatial inspection of the saved DenseNet121 by using gradients of a target class score with respect to its final convolutional representation. The function \texttt{src/gradcam.py} computes the maps from the checkpoint that produced Table~\ref{tab:final} (its SHA-256 is stored with the run), with the network in evaluation mode.

For a binary output logit $z$, define $s_1=z$ for PNEUMONIA and $s_0=-z$ for NORMAL. Let $A^k$ denote feature channel $k$ after the final DenseNet feature activation, of spatial size $u\times v=16\times16$ for a $512\times512$ input. The attribution is
\[
\alpha_k^{(c)}=\frac{1}{uv}\sum_{i,j}\frac{\partial s_c}{\partial A^k_{ij}},
\qquad
L^{(c)}=\operatorname{ReLU}\!\left(\sum_k\alpha_k^{(c)}A^k\right).
\]
Maps are bilinearly enlarged to the input grid. A positive map indicates spatial support for the chosen class score, rather than all features used by the model. The original predicted class is selected as the target and held fixed for the corresponding padding-control map, including if the class prediction changes. Differentiating the logit avoids using the potentially saturated sigmoid probability as the explanation target.

\subsection{Original input and a controlled padding perturbation}
The paired comparison uses \textbf{the same checkpoint in both columns}. The original input follows the deterministic validation/test transform: grayscale decoding replicated into three channels, longest-side resizing to 512 pixels with area interpolation, centered padding, and per-channel normalization with mean 0.482 and standard deviation 0.236. The exact resized image footprint defines a binary mask $M$ of synthetic padding; it does not identify lung tissue.

For the controlled input, set only pixels in $M$ to zero after normalization, corresponding to a uniform raw intensity of 0.482. All pixels inside the resized image footprint, the aspect ratio, image placement, and network weights remain identical. This is an inference-time padding ablation. It tests sensitivity to reflected border content, while retaining the padding extent and its potential geometric information. It therefore cannot be interpreted as removal of every geometry cue. A uniformly filled region also changes the input distribution, so an output change is evidence of sensitivity to this intervention rather than proof of a specific shortcut mechanism.


\subsection{Case selection and test-wide measurement}
For the figures, the two most confident test cases of each outcome (TN, TP, FP, FN) at the
frozen threshold are selected; no map appearance enters the selection. The padding statistics
are then computed on \emph{all} test images. For a nonzero Grad-CAM map with a nonempty padding
region $M$, the padding attribution mass and its area-adjusted enrichment are
\[
m_{\mathrm{pad}}=\frac{\sum_{i,j} M_{ij}L^{(c)}_{ij}}{\sum_{i,j}L^{(c)}_{ij}},
\qquad
e_{\mathrm{pad}}=\frac{m_{\mathrm{pad}}}{|M|/(512^2)}.
\]
An enrichment above one means more positive attribution per pixel in padding than a spatially
uniform map would assign; it is not a causal contribution.


\begin{figure}[htbp]
\centering
\includegraphics[width=\linewidth]{results/figures/error_gallery.pdf}
\caption{Most confident test cases of each outcome at the frozen threshold
(two per outcome, deterministic selection by confidence).}
\label{fig:errors}
\end{figure}

\begin{figure}[htbp]
\centering
\includegraphics[width=0.9\linewidth]{results/figures/gradcam_cases.pdf}
\caption{Grad-CAM for the predicted class on the original input and on the input with synthetic
padding set to the mean intensity (same checkpoint, same target class). Cyan: resized image
footprint, not a lung outline. Colours share one maximum within each row.}
\label{fig:gradcam}
\end{figure}

\begin{figure}[htbp]
\centering
\includegraphics[width=\linewidth]{results/figures/padding_summary.pdf}
\caption{Padding analysis on all test images. Left: share of Grad-CAM mass inside padding
versus padding area (points above the diagonal concentrate attribution in padding). Right:
PNEUMONIA probability before and after neutralising padding.}
\label{fig:padding}
\end{figure}

\fillin{Report the median padding mass vs median padding area, the share of images with
enrichment $>1$ and the share of predictions flipped by neutralisation, for the reference and
for the padding-ablated model (Q3). This turns the four anecdotal cases of the first revision
into a cohort-level measurement.}

\section{Limitations}
\label{sec:limitations}
\begin{itemize}
\item Single-centre pediatric data; no external test set. Results do not transfer to adults
or other hospitals without new evidence.
\item Patient keys are inferred from filenames; a documented identifier would be needed to
certify patient independence.
\item Image geometry separates the classes; the comparison with the geometry baseline and the
padding analysis bound, but do not eliminate, the use of acquisition cues.
\item Grad-CAM has $16\times16$ resolution and no lung masks are available, so attribution
inside the image cannot be equated with pulmonary evidence.
\end{itemize}


\section{Conclusion}
The dataset audit shows that native image geometry alone reaches test ROC-AUC \GeomTestAUC{}, so this value, rather than 0.5, is the reference a neural network must exceed. With a patient-grouped split and validation-only selection, the final configuration (\FinalName) reaches test ROC-AUC \FinalTestAUC{} over \FinalSeeds{} seeds, with sensitivity \FinalTestSens{} and specificity \FinalTestSpec{} at a threshold frozen on validation.
\fillin{One sentence each: effect of leakage (Q0); value of transfer learning and choice of DenseNet121 (Q1); retained imbalance strategy (Q2); effect of the padding ablation on performance and on attribution (Q3).}
Patient identity inferred from filenames, the single-centre source and the coarse resolution of Grad-CAM remain the main limitations.

\appendix
% =============================================================================
% Compressed replacement for the historical results (current Section 8).
% Place after the bibliography or before it, preceded by \appendix.
% =============================================================================
\section{Results of the First Submission}
\label{app:previous}
For traceability, Table~\ref{tab:previous} reports the run of the first submission
(image-level split, seven epochs, $512\times512$). Its thresholds were compared on the test set,
so none of them is an independently selected operating point; its checkpoint was not preserved
with a hash and could not be reproduced. These two problems motivated the protocol of
Section~\ref{sec:protocol}.

\begin{table}[htbp]
\centering
\caption{First submission, source test set (ROC-AUC 0.9607). Retrospective threshold comparison.}
\label{tab:previous}
\small
\begin{tabular}{rrrrrrr}
\toprule
Threshold & TN & FP & FN & TP & Sensitivity & Specificity \\
\midrule
0.40 & 139 & 95 & 2 & 388 & 0.995 & 0.594 \\
0.50 & 165 & 69 & 5 & 385 & 0.987 & 0.705 \\
0.70 & 201 & 33 & 11 & 379 & 0.972 & 0.859 \\
0.80 & 208 & 26 & 26 & 364 & 0.933 & 0.889 \\
\bottomrule
\end{tabular}
\end{table}


\clearpage
\begin{thebibliography}{99}
\bibitem{huang2017densenet}
G. Huang, Z. Liu, L. van der Maaten, and K. Q. Weinberger.
\textit{Densely Connected Convolutional Networks}.
Proceedings of the IEEE Conference on Computer Vision and Pattern Recognition (CVPR), 2017.
\href{https://arxiv.org/abs/1608.06993}{arXiv:1608.06993}.

\bibitem{rajpurkar2017chexnet}
P. Rajpurkar, J. Irvin, K. Zhu, et al.
\textit{CheXNet: Radiologist-Level Pneumonia Detection on Chest X-Rays with Deep Learning}.
2017.
\href{https://arxiv.org/abs/1711.05225}{arXiv:1711.05225}.

\bibitem{kermany2018}
D. S. Kermany, M. Goldbaum, W. Cai, et al.
\textit{Identifying Medical Diagnoses and Treatable Diseases by Image-Based Deep Learning}.
Cell 172(5), 1122--1131.e9, 2018.
\href{https://doi.org/10.1016/j.cell.2018.02.010}{doi:10.1016/j.cell.2018.02.010}.

\bibitem{kermanydata}
D. Kermany, K. Zhang, and M. Goldbaum.
\textit{Labeled Optical Coherence Tomography (OCT) and Chest X-Ray Images for Classification}.
Mendeley Data, version 2, 2018. CC BY 4.0.
\href{https://doi.org/10.17632/rscbjbr9sj.2}{doi:10.17632/rscbjbr9sj.2}.
Chest-radiograph archive: \texttt{ChestXRay2017.zip}.

\bibitem{lin2017focal}
T.-Y. Lin, P. Goyal, R. Girshick, K. He, and P. Doll\'{a}r.
\textit{Focal Loss for Dense Object Detection}.
Proceedings of the IEEE International Conference on Computer Vision (ICCV), 2017.
\href{https://arxiv.org/abs/1708.02002}{arXiv:1708.02002}.

\bibitem{zech2018}
J. R. Zech, M. A. Badgeley, M. Liu, A. B. Costa, J. J. Titano, and E. K. Oermann.
\textit{Variable generalization performance of a deep learning model to detect pneumonia in chest radiographs: A cross-sectional study}.
PLoS Medicine 15(11), e1002683, 2018.
\href{https://doi.org/10.1371/journal.pmed.1002683}{doi:10.1371/journal.pmed.1002683}.

\bibitem{saito2015pr}
T. Saito and M. Rehmsmeier.
\textit{The Precision-Recall Plot Is More Informative than the ROC Plot When Evaluating Binary Classifiers on Imbalanced Datasets}.
PLoS ONE 10(3), e0118432, 2015.
\href{https://doi.org/10.1371/journal.pone.0118432}{doi:10.1371/journal.pone.0118432}.

\bibitem{vickers2006decision}
A. J. Vickers and E. B. Elkin.
\textit{Decision Curve Analysis: A Novel Method for Evaluating Prediction Models}.
Medical Decision Making 26(6), 565--574, 2006.
\href{https://doi.org/10.1177/0272989X06295361}{doi:10.1177/0272989X06295361}.

\bibitem{sklearnthreshold}
scikit-learn developers.
\textit{Tuning the decision threshold for class prediction}.
Official documentation, accessed October 26, 2026.
\url{https://scikit-learn.org/stable/modules/classification_threshold.html}.

\bibitem{selvaraju2017gradcam}
R. R. Selvaraju, M. Cogswell, A. Das, R. Vedantam, D. Parikh, and D. Batra.
\textit{Grad-CAM: Visual Explanations from Deep Networks via Gradient-Based Localization}.
Proceedings of the IEEE International Conference on Computer Vision (ICCV), 618--626, 2017.
\href{https://openaccess.thecvf.com/content_ICCV_2017/papers/Selvaraju_Grad-CAM_Visual_Explanations_ICCV_2017_paper.pdf}{CVF Open Access}.

\bibitem{adebayo2018sanity}
J. Adebayo, J. Gilmer, M. Muelly, I. Goodfellow, M. Hardt, and B. Kim.
\textit{Sanity Checks for Saliency Maps}.
Advances in Neural Information Processing Systems 31, 2018.
\href{https://proceedings.neurips.cc/paper/2018/hash/294a8ed24b1ad22ec2e7efea049b8737-Abstract.html}{NeurIPS proceedings}.

\end{thebibliography}

\end{document}